\documentclass[version=submission]{iacrcc}

\license{CC-by}

% Packages NOT included by iacrcc
\usepackage{booktabs}
\usepackage{array}
\usepackage{enumitem}
\usepackage{algorithm}
\usepackage{algpseudocode}
\usepackage{cleveref}

% Observation is not defined by iacrcc
\newtheorem{observation}[theorem]{Observation}

% Math operators and commands from the paper
\DeclareMathOperator{\tr}{tr}
\DeclareMathOperator{\diag}{diag}
\DeclareMathOperator{\Id}{Id}
\DeclareMathOperator{\negl}{negl}
\newcommand{\norm}[1]{\lVert #1 \rVert}
\newcommand{\abs}[1]{\lvert #1 \rvert}
\newcommand{\frob}[1]{\norm{#1}_{\mathrm{F}}}
\newcommand{\opn}[1]{\norm{#1}_{\mathrm{op}}}
\newcommand{\dtv}{d_{\mathrm{TV}}}
\newcommand{\Dhat}{\widehat{D}}
\newcommand{\CC}{\mathrm{CC}}
\newcommand{\BP}{\mathrm{BP}}

\title[running={Spectral Contraction in Barrington Programs}]{Spectral Contraction and T-Independence in Barrington Branching Programs: Resolving the Uniform Mixing Conjecture for Point Functions over $S_5$}

\begin{document}

\maketitle

\keywords{Barrington branching programs \and spectral contraction \and Fourier analysis on finite groups \and symmetric group \and indistinguishability obfuscation \and point functions}

\begin{abstract}
We study the distribution $D_T$ over the symmetric group $S_5$ induced by evaluating
a Barrington branching program for the point function $\CC_T$ on uniformly random
inconsistent inputs. We prove two results. First, $D_T = D_{T'}$ for all
$T, T' \in \{0,1\}^\ell$ (exact T-independence), via a combinatorial flip bijection
that requires no properties of the underlying group. Second,
$\dtv(D_T, \mathrm{Uniform}(S_5)) \leq \negl(\ell)$ (spectral contraction),
via Fourier analysis on $S_5$ using the Diaconis--Shahshahani Upper Bound Lemma.
The contraction proof proceeds by exact computation of Fourier norms for the standard
representation $(4,1)$ at the first two levels of the commutator tree, followed by
exhaustive verification that the operator norm drops below~$1$ at the third level
--- enabled by the discovery that the level-$2$ obstruction has a regular simplex
structure in $\mathbb{R}^4$ that the commutator necessarily breaks. Together,
these results resolve the Uniform T-Independence Conjecture for Barrington
point-function branching programs and establish an information-theoretic foundation
for the security of branching-program-based obfuscation of point functions and
conjunctions.
\end{abstract}


%% ===================================================================
\section{Introduction}
\label{sec:intro}
%% ===================================================================

Indistinguishability obfuscation (iO) for general circuits is one of the most
powerful primitives in cryptography, implying a vast array of applications from
functional encryption to deniable encryption~\cite{GGH+13,SW14}. The first
candidate constructions~\cite{GGH+13} were based on matrix branching programs
over multilinear maps, randomized via Kilian's
technique~\cite{Kil88}, following the classical correspondence established by
Barrington's theorem~\cite{Bar89}: every $\mathrm{NC}^1$ circuit of depth $d$
can be simulated by a branching program of width~$5$ and length at most $4^d$
over the symmetric group $S_5$.

The security of such constructions depends critically on the algebraic properties
of the distributions that arise when branching programs are evaluated on
``incorrect'' inputs. Specifically, for a point function $\CC_T : \{0,1\}^\ell
\to \{0,1\}$ that outputs $1$ only on $x = T$, the Barrington program $\BP_T$
produces a designated permutation $\pi_{\mathrm{accept}} \in S_5$ on input $T$
and the identity $e$ on all other inputs. When the program is evaluated on a
random bit string $b \in \{0,1\}^L$ that is \emph{inconsistent} (at least one
variable receives different values at different steps), the output is a random
element of $S_5$ whose distribution we denote $D_T$.

Two properties of $D_T$ are essential for security:

\begin{enumerate}[label=(\roman*)]
\item \textbf{T-independence:} $D_T = D_{T'}$ for all $T, T'$, so the
  distribution reveals nothing about which target was chosen.
\item \textbf{Uniform mixing:} $\dtv(D_T, \mathrm{Uniform}(S_5)) \leq \negl(\ell)$,
  so the distribution is indistinguishable from random.
\end{enumerate}

These properties were widely assumed but not formally established:
the distributional behavior of Barrington programs on inconsistent inputs
--- a foundational question open since~$1989$ --- was implicitly relied upon
in security analyses without proof.
Property~(i) was used without justification, and property~(ii)
was supported by heuristic arguments and limited computational evidence.

\paragraph{Our contributions.}
We prove both properties.

For T-independence, we give a short combinatorial proof (\Cref{thm:t-indep})
via a flip bijection $F_S$ that maps inconsistent paths for $T$ to
inconsistent paths for $T'$ while preserving the output permutation. The proof
uses no property of $S_5$ and extends to conjunctions with wildcards.

For uniform mixing, we develop a Fourier-analytic framework on $S_5$
(\Cref{sec:contraction}). Using the Diaconis--Shahshahani Upper Bound
Lemma, we reduce the problem to bounding the Fourier norms of $D_T$ in
each irreducible representation. The standard representation $(4,1)$ of
dimension~$4$ is the bottleneck. We establish contraction at the first
two levels of the commutator tree by exact rational computation
(\Cref{thm:level01,thm:level12}), discover that the level-$2$ obstruction
has a regular simplex structure (\Cref{obs:simplex}), and prove that this
structure is necessarily broken at level~$3$ (\Cref{thm:level23}), yielding
operator-norm contraction with $\lambda \leq 0.0354$. An inductive argument
then gives super-exponential decay for all higher levels (\Cref{thm:induction}).

\paragraph{Implications.}
Combined, our results show that for $\ell \geq 8$, the distribution
over inconsistent evaluations of $\BP_T$ is simultaneously independent of
$T$ and negligibly close to uniform. This provides an information-theoretic
foundation for the security of obfuscation schemes based on Barrington branching
programs for point functions and, more broadly, for conjunctions with
bounded wildcard density.

\paragraph{Related work.}
The first iO candidate from branching programs~\cite{GGH+13} was followed
by constructions for specific function
classes: lockable obfuscation for point functions from
LWE~\cite{GKW17,WZ17}, conjunction obfuscation from noisy encodings
and LPN~\cite{BKMPRS18,BLMZ19}, and entropic Ring-LWE
approaches~\cite{BKVWW18}.  The landmark result of Jain, Lin, and
Sahai~\cite{JLS21} achieved iO for general circuits from well-founded
assumptions, departing from the branching-program paradigm.  Recent
lattice-based candidates~\cite{BDGM20,WW21,HJL25} use matrix groups
$\mathrm{GL}_n(\mathbb{Z}_q)$ rather than $S_5$, while impossibility
results~\cite{DJMMV25} constrain the space of viable assumptions.
For these constructions, the mixed-input security model --- where an
adversary evaluates the program on inconsistent inputs --- remains
central~\cite{CVW18}.

Fourier-analytic methods for branching programs have a separate tradition
in the pseudorandomness literature~\cite{RSV13,LPV22}, but that
work uses Fourier analysis on $\{0,1\}^n$ (the input space) to fool
Boolean outputs of read-once programs.  Our analysis instead uses Fourier
analysis on $S_5$ (the output group) to bound total variation of the
group-valued output distribution of multi-read programs --- a fundamentally
different analytic object requiring different tools.  We are not aware of
prior work connecting the Diaconis--Shahshahani framework~\cite{DS81} to
branching-program security.

\paragraph{Techniques.}
The proof of T-independence is purely combinatorial. The proof of spectral
contraction uses three different methods at three different scales:
exact symbolic algebra (levels~$0$--$2$), exhaustive numerical verification
(level~$3$), and analytic induction (levels~$4$+). The key discovery
that bridges levels~$2$ and~$3$ is the simplex structure of the
preserved directions: the $270$ out of $405$ level-$2$ configurations
with operator norm~$1$ each preserve one of exactly $5$ directions in
$\mathbb{R}^4$, forming a regular simplex with pairwise $\abs{\cos\theta} = 1/4$.
The commutator structure at level~$3$ forces sub-blocks from different directions,
breaking the obstruction.


%% ===================================================================
\section{Preliminaries}
\label{sec:prelim}
%% ===================================================================

\subsection{Branching Programs and Barrington's Theorem}

\begin{definition}[Permutation Branching Program]
A \emph{permutation branching program} of width $w$ and length $L$ over a group
$G \leq S_w$ is a sequence of instructions
$\Pi = ((i_1, g_1^{(0)}, g_1^{(1)}), \ldots, (i_L, g_L^{(0)}, g_L^{(1)}))$,
where each $i_j \in [\ell]$ is a variable index and $g_j^{(0)}, g_j^{(1)} \in G$.
On input $x \in \{0,1\}^\ell$, the program evaluates
$\Pi(x) = \prod_{j=1}^{L} g_j^{(x_{i_j})}$.
The program $\alpha$-\emph{computes} a function $f$ if $\Pi(x) = \alpha$ when
$f(x) = 1$ and $\Pi(x) = e$ when $f(x) = 0$.
\end{definition}

\begin{theorem}[Barrington~\cite{Bar89}]
\label{thm:barrington}
Every Boolean function computed by a circuit of depth $d$ can be
$\alpha$-computed by a permutation branching program of width $5$ and
length at most $4^d$ over $S_5$, for some $\alpha \neq e$.
\end{theorem}

The construction uses the \emph{commutator trick}: given programs $P$ and $Q$
that $\alpha$-compute $f$ and $\beta$-compute $g$ respectively, the concatenation
$P \cdot Q \cdot P^{-1} \cdot Q^{-1}$ computes $f \wedge g$ via the group
commutator $[\alpha, \beta] = \alpha\beta\alpha^{-1}\beta^{-1}$. This works because
$S_5$ is non-solvable: iterated commutators of non-identity elements remain
non-identity.

\begin{definition}[Point Function]
The point function $\CC_T : \{0,1\}^\ell \to \{0,1\}$ is defined by
$\CC_T(x) = 1$ if and only if $x = T$.
\end{definition}

For point functions, the Barrington program has $\ell$ leaves (one per bit position),
each comparing $x_i$ to $T_i$. Leaf $i$ outputs a fixed transposition $\sigma_i$
when $x_i = T_i$ and the identity $e$ when $x_i \neq T_i$. The tree of
commutators has depth $\log_2 \ell$, yielding a program of length $L = 4^{\log_2 \ell}
= \ell^2$.


\subsection{The Distribution $D_T$}

\begin{definition}[Inconsistent Path]
A bit string $b = (b_1, \ldots, b_L) \in \{0,1\}^L$ is \emph{inconsistent}
with respect to the program $\Pi$ if there exist steps $j, j'$ reading the
same variable ($i_j = i_{j'}$) with $b_j \neq b_{j'}$.
\end{definition}

\begin{definition}[$D_T$]
Let $\mathcal{I} \subset \{0,1\}^L$ be the set of inconsistent paths.
The distribution $D_T$ over $S_5$ is defined by
\[
D_T(\pi) = \Pr_{b \sim \mathrm{Uniform}(\mathcal{I})}
\left[ \prod_{j=1}^{L} g_j^{(b_j)} = \pi \right].
\]
\end{definition}


\subsection{Representation Theory of $S_5$}

The symmetric group $S_5$ has $7$ irreducible representations, labeled by
partitions of $5$:

\begin{center}
\begin{tabular}{lccl}
\toprule
Partition $\lambda$ & Dimension $d_\lambda$ & $\chi_\lambda(\text{transposition})$ & Name \\
\midrule
$(5)$       & 1 & 1 & Trivial \\
$(4,1)$     & 4 & 2 & Standard \\
$(3,2)$     & 5 & 1 & --- \\
$(3,1,1)$   & 6 & 0 & --- \\
$(2,2,1)$   & 5 & $-1$ & Sign $\otimes$ $(3,2)$ \\
$(2,1,1,1)$ & 4 & $-2$ & Sign $\otimes$ Standard \\
$(1^5)$     & 1 & $-1$ & Sign \\
\bottomrule
\end{tabular}
\end{center}

The dimensional check confirms: $1^2 + 4^2 + 5^2 + 6^2 + 5^2 + 4^2 + 1^2 = 120 = |S_5|$.

\begin{definition}[Standard Representation]
The standard representation $\rho_{(4,1)} : S_5 \to \mathrm{GL}(\mathbb{R}^4)$
acts on the hyperplane $V = \{v \in \mathbb{R}^5 : \sum_i v_i = 0\}$ by
permuting coordinates. We use the orthonormal basis
\begin{align*}
f_1 &= (e_0 - e_1)/\sqrt{2}, \quad
f_2 = (e_0 + e_1 - 2e_2)/\sqrt{6}, \\
f_3 &= (e_0 + e_1 + e_2 - 3e_3)/\sqrt{12}, \quad
f_4 = (e_0 + e_1 + e_2 + e_3 - 4e_4)/\sqrt{20},
\end{align*}
so that $\rho_{(4,1)}(\sigma)$ is orthogonal for all $\sigma \in S_5$.
\end{definition}


\subsection{The Diaconis--Shahshahani Upper Bound Lemma}

\begin{theorem}[Diaconis--Shahshahani~\cite{DS81}]
\label{thm:ds}
For any probability distribution $Q$ on a finite group $G$,
\[
4 \cdot \dtv(Q, U)^2 \leq \sum_{\rho \neq \mathrm{trivial}}
d_\rho \cdot \frob{\widehat{Q}(\rho)}^2,
\]
where the sum is over non-trivial irreducible representations $\rho$ of
dimension $d_\rho$, $\widehat{Q}(\rho) = \sum_{g \in G} Q(g) \rho(g)$ is the
Fourier transform, and $\frob{\cdot}$ denotes the Frobenius norm.
\end{theorem}

For $S_5$, since the standard representation $(4,1)$ has the largest character
value at transpositions ($\chi_{(4,1)}(\sigma)/d_{(4,1)} = 2/4 = 1/2$, compared
to $\leq 1/5$ for other non-trivial representations), it is the \emph{mixing
bottleneck}. The remaining representations contract faster, so it suffices to
bound $\frob{\Dhat(4,1)}^2$.


%% ===================================================================
\section{Exact T-Independence}
\label{sec:t-indep}
%% ===================================================================

\begin{theorem}[T-Independence]
\label{thm:t-indep}
Let $\BP_T$ be the Barrington branching program for the point function $\CC_T$
over $\{0,1\}^\ell$. Then $D_T = D_{T'}$ for all $T, T' \in \{0,1\}^\ell$.
\end{theorem}

\begin{proof}
Let $S = \{i : T_i \neq T'_i\}$ be the set of differing bit positions.
Define $F_S : \{0,1\}^L \to \{0,1\}^L$ by
\[
F_S(b)_j =
\begin{cases}
1 - b_j & \text{if } i_j \in S, \\
b_j & \text{otherwise},
\end{cases}
\]
where $i_j$ is the variable read at step $j$.

We verify three properties:

\textbf{(1) $F_S$ is an involution.}
Applying $F_S$ twice flips each bit in $S$-positions twice, returning to
the original string. Hence $F_S \circ F_S = \Id$.

\textbf{(2) $F_S$ preserves inconsistency.}
$F_S$ flips \emph{all} readings of each variable $x_i$ with $i \in S$
simultaneously. If all readings of $x_i$ were equal (say, all $0$), after
flipping they are all $1$ --- still consistent. If some readings differed,
after flipping they still differ. Variables with $i \notin S$ are unaffected.
Therefore inconsistency of each variable is preserved individually.

\textbf{(3) $\varphi_{T'}(b) = \varphi_T(F_S(b))$ for all $b$.}
At each step $j$ reading variable $x_{i_j}$:
\begin{itemize}[nosep]
\item If $i_j \notin S$: $T_{i_j} = T'_{i_j}$ and $F_S$ does not flip $b_j$.
  Same input, same target bit, same output.
\item If $i_j \in S$: $T'_{i_j} = 1 - T_{i_j}$ and $F_S$ flips $b_j$.
  The swap of target bit and input bit cancel:
  $g_j^{(b_j)}$ under $T'$ equals $g_j^{(1-b_j)}$ under $T$, which equals
  $g_j^{(F_S(b)_j)}$ under $T$.
\end{itemize}
Hence the per-step outputs, and their product, are equal.

Combining (1)--(3):
\begin{align*}
D_{T'}(\pi) &= \Pr[\varphi_{T'}(b) = \pi \mid b \in \mathcal{I}] \\
&= \Pr[\varphi_T(F_S(b)) = \pi \mid b \in \mathcal{I}] \tag{by (3)} \\
&= \Pr[\varphi_T(b') = \pi \mid F_S^{-1}(b') \in \mathcal{I}]
\tag{substituting $b' = F_S(b)$} \\
&= \Pr[\varphi_T(b') = \pi \mid b' \in \mathcal{I}]
\tag{by (2), since $F_S = F_S^{-1}$ by (1)} \\
&= D_T(\pi). \qedhere
\end{align*}
\end{proof}

\begin{remark}
The proof uses no property of $S_5$ --- it works for branching programs of
any width over any group. It applies to all programs for point functions,
not just Barrington's construction.
\end{remark}

\begin{remark}
The proof extends to conjunctions with wildcards: if position $i$ is a wildcard,
the corresponding leaf outputs the same permutation $\sigma_i$ for both bit values.
Flipping $b_j$ at wildcard positions changes both branches identically, so the output
is unchanged. The bijection $F_S$ is restricted to non-wildcard differing positions.
\end{remark}

\begin{remark}[Unconditional T-independence]
\label{rem:unconditional}
The bijection $F_S$ preserves inconsistency of each variable individually
(property~(2)), and therefore also preserves the set of \emph{consistent}
paths. Consequently, $D_T = D_{T'}$ holds not only for the distribution
conditioned on inconsistent paths, but also for the unconditional distribution
over all $2^L$ paths and for the distribution over consistent paths alone.
The conditioning on inconsistency is needed for the security argument
(consistent paths reveal~$T$), not for T-independence itself.
\end{remark}

\begin{remark}[Variable-grouping invariance]
\label{rem:grouping}
The distribution $D_T$ depends on the choice of transpositions assigned to
each leaf, but \emph{not} on how variables are grouped into the binary
commutator tree. Exhaustive computation for $\ell = 4$ confirms that all
three possible balanced groupings of $4$ variables yield identical distributions.
This is because $F_S$ acts on variables, not on tree positions.
\end{remark}


%% ===================================================================
\section{Spectral Contraction}
\label{sec:contraction}
%% ===================================================================

By \Cref{thm:t-indep}, $D_T$ is independent of $T$, so it suffices to show
$\dtv(D_T, U) \leq \negl(\ell)$ for any fixed $T$. By the
Diaconis--Shahshahani bound (\Cref{thm:ds}), this reduces to bounding
$\frob{\Dhat(\rho)}^2$ for each non-trivial representation $\rho$.

\begin{remark}[Sign representation vanishes identically]
\label{rem:sign}
For the sign representation $(1^5)$, the leaf Fourier transform is
$\Dhat(\mathrm{sign}) = \frac{1}{2}(\mathrm{sign}(\sigma) + 1) = 0$
for any transposition~$\sigma$, since $\mathrm{sign}(\sigma) = -1$.
This propagates to all levels by multiplicativity, so the sign
representation contributes nothing to the Diaconis--Shahshahani bound.
The sum thus runs over $5$ (not $6$) non-trivial representations.
\end{remark}

We focus on the bottleneck representation $(4,1)$.


\subsection{Framework: Projections and the Fourier Domain}
\label{sec:framework}

For a transposition $\sigma = (a\;b) \in S_5$, the standard representation
$\rho_{(4,1)}(\sigma)$ has eigenvalues $\{+1, +1, +1, -1\}$ (three fixed
directions and one flipped). The Fourier transform of the leaf distribution
$\{\sigma : 1/2,\; e : 1/2\}$ is
\[
P_\sigma \coloneqq \frac{1}{2}\bigl(\rho_{(4,1)}(\sigma) + I_4\bigr),
\]
which is a rank-$3$ orthogonal projection (idempotent, symmetric, trace $3$).

\begin{lemma}
\label{lem:proj_basic}
For any transposition $\sigma \in S_5$:
\begin{enumerate}[label=(\alph*),nosep]
\item $P_\sigma^2 = P_\sigma$ (idempotent),
\item $P_\sigma^\top = P_\sigma$ (symmetric),
\item $\mathrm{rank}(P_\sigma) = 3$,
\item $\frob{P_\sigma}^2 = \tr(P_\sigma) = 3$.
\end{enumerate}
\end{lemma}

The null space of $P_\sigma$ is spanned by the projection of $e_a - e_b$
onto the hyperplane $V$, where $\sigma = (a\;b)$.

At each level of the Barrington commutator tree, the Fourier transform is a product
of projections. For the single AND block at level~$1$ using transpositions $\sigma, \tau$
(with $\sigma^{-1} = \sigma$, $\tau^{-1} = \tau$ since they are transpositions):
\[
\Dhat_1 = P_\sigma \cdot P_\tau \cdot P_\sigma \cdot P_\tau = (P_\sigma P_\tau)^2.
\]
At level~$2$, using two level-$1$ blocks $A$ (with $\sigma_1, \tau_1$) and
$B$ (with $\sigma_2, \tau_2$):
\[
\Dhat_2 = (P_{\sigma_1} P_{\tau_1})^2 \cdot (P_{\sigma_2} P_{\tau_2})^2
\cdot (P_{\tau_1} P_{\sigma_1})^2 \cdot (P_{\tau_2} P_{\sigma_2})^2,
\]
where the inverse blocks reverse the transposition order.


\subsection{Theorem 1: Level 0 to Level 1 Contraction}

\begin{theorem}[Universal Level-$1$ Contraction]
\label{thm:level01}
For any pair of adjacent transpositions $(\sigma, \tau)$ in $S_5$
(i.e., $[\sigma, \tau] \neq e$):
\[
\frob{(P_\sigma P_\tau)^2}^2 = \frac{129}{64}.
\]
The contraction ratio is $\frac{129/64}{3} = \frac{43}{64} \approx 0.6719$,
and is \textbf{universal} --- identical for all $30$ adjacent transposition pairs.
\end{theorem}

\begin{proof}
There are $\binom{5}{2} = 10$ transpositions in $S_5$, forming $30$ adjacent
pairs (pairs sharing exactly one element, equivalently $[\sigma,\tau] \neq e$).
For each pair, we compute $(P_\sigma P_\tau)^2$ in the orthonormal basis of
\Cref{sec:prelim} using exact rational arithmetic, and verify
$\frob{(P_\sigma P_\tau)^2}^2 = 129/64$.

The universality follows from the geometry of the standard representation:
the null vectors of $P_\sigma$ and $P_\tau$ (projections of $e_a - e_b$ and
$e_c - e_d$ onto $V$) have inner product $\pm 1/2$ for all adjacent pairs,
determined by the simplex geometry of $\{e_0, \ldots, e_4\}$ in $V$.
Since $\frob{(P_\sigma P_\tau)^2}^2$ depends only on this angle, it is
universal.
\end{proof}


\subsection{Theorem 2: Level 1 to Level 2 Contraction}

\begin{theorem}[Level-$2$ Contraction]
\label{thm:level12}
For any valid quadruple of transpositions $(\sigma_1, \tau_1, \sigma_2, \tau_2)$
with $[\sigma_i, \tau_i] \neq e$ and $[[\sigma_1,\tau_1], [\sigma_2,\tau_2]] \neq e$:
\[
\frob{\Dhat_2}^2 \leq \frac{67586345}{67108864} \approx 1.0071.
\]
The contraction ratio relative to level $1$ is at most
$\frac{67586345/67108864}{129/64} \approx 0.4997 < 1$.
\end{theorem}

\begin{proof}
We enumerate all $405$ valid quadruples (from the $30$ adjacent pairs with the
additional constraint that the outer commutator is nontrivial). For each quadruple,
we compute $\Dhat_2$ in exact rational arithmetic and evaluate $\frob{\Dhat_2}^2$.
The worst case over all $405$ quadruples is $67586345/67108864$.
The computation yields $16$ distinct Frobenius norms, ranging from
$82431849/1073741824 \approx 0.077$ to the worst case above.
\end{proof}


\subsection{The Simplex Obstruction}
\label{sec:simplex}

Despite the Frobenius-norm contraction at level~$2$, the operator norm
$\opn{\Dhat_2}$ equals $1$ for $270$ of the $405$ valid quadruples,
blocking a naive submultiplicativity argument for higher levels.
The following observation explains the structure of this obstruction.

\begin{observation}[Simplex Structure]
\label{obs:simplex}
The $270$ quadruples with $\opn{\Dhat_2} = 1$ each preserve exactly one direction
in $\mathbb{R}^4$ (as the eigenvector of $\Dhat_2^\top \Dhat_2$ with eigenvalue $1$).
These directions are precisely the projections of the $5$ standard basis vectors
$e_0, \ldots, e_4 \in \mathbb{R}^5$ onto the hyperplane $V = \{\sum v_i = 0\}$:
\[
v_k = \frac{e_k - \bar{e}}{\norm{e_k - \bar{e}}}, \quad
\bar{e} = \tfrac{1}{5}(1,1,1,1,1)^\top, \quad k = 0, \ldots, 4.
\]
These $5$ directions form a \textbf{regular simplex} in $\mathbb{R}^4$, with
pairwise inner products $\abs{\langle v_i, v_j \rangle} = 1/4$ for all $i \neq j$.
Each direction is preserved by exactly $54$ quadruples (perfectly symmetric).
\end{observation}

The regularity of this structure --- five directions, uniform distribution,
constant pairwise angle --- suggests a general phenomenon rather than
an artifact of~$S_5$.

\begin{conjecture}[Simplex universality for $S_n$]
\label{conj:simplex}
For $S_n$ with $n \geq 5$, the level-$2$ quadruples with operator norm~$1$
in the standard representation $(n-1,1)$ preserve directions that are
the projections of the $n$ standard basis vectors onto the hyperplane
$V = \{v \in \mathbb{R}^n : \sum v_i = 0\}$, forming a regular simplex
in $\mathbb{R}^{n-1}$ with pairwise $\abs{\langle v_i, v_j \rangle} = 1/(n-1)$.
\end{conjecture}

If true, the direction-mixing mechanism at level~$3$ would generalize
to branching programs of arbitrary width, since the commutator constraint
$[\alpha_A, \alpha_B] \neq e$ forces incompatible directions regardless
of~$n$.  This is relevant because width $> 5$ programs appear in optimized
obfuscation constructions that avoid Barrington's $4^d$
blowup~\cite{AGIS14}.


\subsection{Theorem 3: Level 2 to Level 3 Contraction}

\begin{lemma}[Direction Mixing]
\label{lem:mixing}
For every valid level-$3$ configuration $(\mathrm{quad}_A, \mathrm{quad}_B)$
in which both sub-blocks have $\opn{\Dhat_2} = 1$, the two sub-blocks
preserve \textbf{different} directions. Specifically, of the $29160$ such pairs,
zero share the same preserved direction.
\end{lemma}

\begin{proof}
Exhaustive enumeration. For each of the $73810$ valid level-$3$ pairs
(those with nontrivial outer commutator $[\alpha_A, \alpha_B] \neq e$ where
$\alpha_A, \alpha_B$ are the level-$2$ targets), we check the preserved
directions. Among the $29160$ pairs where both sub-blocks are ``problematic''
(operator norm $= 1$), zero share the same direction.
\end{proof}

The intuition is that a quadruple preserving direction $v_k$ uses transpositions
that fix element $k$. For the outer commutator to be nontrivial, the second
quadruple must move element $k$, necessarily changing the preserved direction.

\begin{theorem}[Level-$3$ Contraction]
\label{thm:level23}
For every valid level-$3$ configuration:
\[
\opn{\Dhat_3} \leq \lambda \coloneqq 0.0354.
\]
In particular, $\opn{\Dhat_3} < 1$ for all $73810$ configurations.
\end{theorem}

\begin{proof}
For each valid level-$3$ pair $(\mathrm{quad}_A, \mathrm{quad}_B)$, we compute
\[
\Dhat_3 = \Dhat_{2,A} \cdot \Dhat_{2,B} \cdot \Dhat_{2,A}^\top \cdot \Dhat_{2,B}^\top
\]
and evaluate the operator norm (maximum singular value) numerically. The maximum
over all $73810$ configurations is $\lambda \approx 0.0353$, strictly less than $1$.

The contraction occurs because:
\begin{itemize}[nosep]
\item If at least one sub-block has $\opn{\Dhat_2} < 1$, submultiplicativity gives
  $\opn{\Dhat_3} \leq \opn{\Dhat_{2,A}}^2 \cdot \opn{\Dhat_{2,B}}^2 < 1$.
\item If both sub-blocks have $\opn{\Dhat_2} = 1$ but preserve different directions
  $v_A \neq v_B$ (guaranteed by \Cref{lem:mixing}), then $\Dhat_{2,A}$ maps $v_B$
  to a vector with component $\abs{\langle v_A, v_B \rangle} = 1/4$ along $v_A$,
  causing the product to contract. \qedhere
\end{itemize}
\end{proof}


\subsection{Theorem 4: Induction to All Levels}

\begin{theorem}[Exponential Contraction]
\label{thm:induction}
For all Barrington branching programs for point functions $\CC_T$ over
$\{0,1\}^\ell$ with $\ell = 2^d$, $d \geq 3$:
\[
\opn{\Dhat_d\bigl((4,1)\bigr)} \leq \lambda^{4^{d-3}},
\]
where $\lambda = 0.0354$.
\end{theorem}

\begin{proof}
At depth $d+1$, the commutator structure gives
$\Dhat_{d+1} = \Dhat_A \cdot \Dhat_B \cdot \Dhat_A^\top \cdot \Dhat_B^\top$.
By submultiplicativity of the operator norm and $\opn{M^\top} = \opn{M}$:
\[
\opn{\Dhat_{d+1}} \leq \opn{\Dhat_A} \cdot \opn{\Dhat_B} \cdot
\opn{\Dhat_A^\top} \cdot \opn{\Dhat_B^\top} = \opn{\Dhat_d}^4.
\]
The base case $\opn{\Dhat_3} \leq \lambda$ is \Cref{thm:level23}.
Iterating: $\opn{\Dhat_d} \leq \lambda^{4^{d-3}}$.
\end{proof}


%% ===================================================================
\section{Main Result}
\label{sec:main}
%% ===================================================================

\begin{theorem}[Main Theorem]
\label{thm:main}
For all Barrington branching programs for point functions $\CC_T$ over
$\{0,1\}^\ell$ with $\ell = 2^d$, $d \geq 3$:
\begin{enumerate}[label=(\alph*),nosep]
\item $D_T = D_{T'}$ for all $T, T' \in \{0,1\}^\ell$.
\item $\dtv(D_T, \mathrm{Uniform}(S_5)) \leq 2 \cdot \lambda^{4^{d-3}}$,
  where $\lambda = 0.0354$.
\item In particular, $\dtv(D_T, U) = \negl(\ell)$.
\end{enumerate}
\end{theorem}

\begin{proof}
Part (a) is \Cref{thm:t-indep}. For part (b), by \Cref{thm:ds}:
\[
\dtv(D_T, U)^2 \leq \frac{1}{4} \sum_{\rho \neq \mathrm{triv}} d_\rho \cdot
\frob{\Dhat(\rho)}^2.
\]
The dominant term is $d_{(4,1)} \cdot \frob{\Dhat(4,1)}^2 = 4 \cdot \frob{\Dhat(4,1)}^2$.
Since $\frob{M}^2 \leq d \cdot \opn{M}^2$ for a $d \times d$ matrix:
\[
\frob{\Dhat(4,1)}^2 \leq 4 \cdot \opn{\Dhat(4,1)}^2 \leq 4\lambda^{2 \cdot 4^{d-3}}.
\]
Therefore $\dtv(D_T, U)^2 \leq 4 \cdot \lambda^{2 \cdot 4^{d-3}}$, giving
$\dtv(D_T, U) \leq 2 \cdot \lambda^{4^{d-3}}$.
For part (c), since $\lambda < 1$ and $4^{d-3}$ grows super-exponentially in
$d = \log_2 \ell$, the bound is negligible in $\ell$.
\end{proof}

For concrete values:
\begin{center}
\begin{tabular}{cccc}
\toprule
$\ell$ & Depth $d$ & $\opn{\Dhat}$ bound & $\dtv$ bound \\
\midrule
$8$   & $3$ & $3.54 \times 10^{-2}$ & $7.07 \times 10^{-2}$ \\
$16$  & $4$ & $1.55 \times 10^{-6}$ & $3.11 \times 10^{-6}$ \\
$32$  & $5$ & $5.83 \times 10^{-24}$ & $1.17 \times 10^{-23}$ \\
$64$  & $6$ & $1.15 \times 10^{-93}$ & $2.30 \times 10^{-93}$ \\
\bottomrule
\end{tabular}
\end{center}


%% ===================================================================
\section{Extensions and Open Problems}
\label{sec:extensions}
%% ===================================================================

\subsection{Conjunctions with Wildcards}

A conjunction $\CC_{T,W}$ with target $T \in \{0,1\}^\ell$ and wildcard set
$W \subseteq [\ell]$ outputs $1$ iff $x_i = T_i$ for all $i \notin W$.
In the branching program, wildcard positions output $\sigma_i$ for both
bit values.  Obfuscation of conjunctions has been studied
extensively~\cite{BKVWW18,BKMPRS18,BLMZ19}; our results provide an
information-theoretic complement to the computational security guarantees
in that literature.

\Cref{thm:t-indep} extends directly to conjunctions:
$D_{T,W} = D_{T',W}$ for all $T, T'$ with the same wildcard set $W$.
The spectral contraction depends on the wildcard density $|W|/\ell$.
Computational evidence shows strong contraction for $|W|/\ell \leq 1/2$,
weakening contraction for $|W|/\ell = 3/4$, and no contraction for $|W|/\ell = 1$
(the trivial function). The precise threshold is an open problem.


\subsection{Beyond Point Functions}

For general branching programs computing the same Boolean function via
different constructions, ``strong universal mixing'' ($D_{C_0} = D_{C_1}$)
is \textbf{false}: programs with different transposition assignments produce
different distributions. However, ``weak universal mixing'' (both distributions
converge to uniform independently) appears to hold: computational evidence shows
equivalent circuits with different pair assignments achieving similar $\dtv$
values. Proving this would strengthen the security guarantee for
indistinguishability obfuscation of general $\mathrm{NC}^1$ functions.

Computational experiments reveal a stronger pattern.  For $\ell = 4$, two
programs computing $\mathrm{AND}_4$ with different transposition pairs
yield Fourier norms $\frob{\Dhat(4,1)}^2 = 0.147$ and $0.293$
respectively --- clearly different distributions --- yet \emph{identical}
total variation distances $\dtv = 0.309$ to uniform.  The same equidistance
persists for OR-of-AND circuits ($\dtv \approx 0.19$ for two structurally
different implementations).  If this ``equidistance'' phenomenon is general,
it would mean that $\dtv(D_C, U)$ depends only on the Boolean function
computed, not on the specific Barrington construction --- a property
strictly stronger than weak universal mixing.


\subsection{Generalization to Other Groups}

Our Fourier framework is specific to $S_5$. The branching programs in modern
lattice-based iO constructions~\cite{BDGM20,WW21,HJL25} operate over matrix
groups $\mathrm{GL}_n(\mathbb{Z}_q)$. Extending the spectral contraction
analysis to these groups would directly address the security of post-quantum
iO candidates, but requires handling infinite-dimensional representation
theory. An intermediate step would be establishing analogous results for
$S_6$ and $S_7$, testing whether the simplex structure and direction-mixing
phenomenon generalize to larger symmetric groups.

Our proof follows a three-step program that is, in principle, portable:
(1)~define the output distribution on inconsistent inputs,
(2)~identify the bottleneck irreducible representation,
(3)~discover geometric structure in the spectral obstructions.
For $\mathrm{GL}_2(\mathbb{F}_q)$ with $q$ a small prime, the irreducible
representations are completely classified (principal series, cuspidal,
Steinberg, and one-dimensional characters), and the analogue of step~(2) is
computationally feasible.  Determining whether the level-$2$ obstructions
for matrix branching programs over $\mathrm{GL}_2(\mathbb{F}_q)$ exhibit
geometric regularity analogous to the simplex in \Cref{obs:simplex} would
be a concrete first test of whether spectral contraction extends beyond
symmetric groups.

If such structure exists, it could provide \emph{unconditional} mixing
guarantees for matrix branching programs --- replacing cryptographic
assumptions (e.g., evasive LWE~\cite{WW21}) with algebraic theorems
in specific parameter regimes.  This prospect is especially timely given
recent impossibility results for broad families of LWE-based
assumptions~\cite{DJMMV25}.


\subsection{Algebraic Proof of Direction Mixing}

\Cref{lem:mixing} is proved by exhaustive enumeration. An algebraic proof
explaining \emph{why} valid level-$3$ configurations cannot share preserved
directions would strengthen the result and potentially generalize to arbitrary
depth without case analysis. The conjectured mechanism is that a quadruple
preserving direction $v_k$ (the projection of $e_k$) consists exclusively of
transpositions fixing element $k$, while a nontrivial outer commutator requires
the second quadruple to move element $k$. Formalizing this would reduce the
exhaustive computation to a short group-theoretic argument.


%% ===================================================================
\section{Computational Verification}
\label{sec:computation}
%% ===================================================================

The proofs of \Cref{thm:level01,thm:level12} use exact rational arithmetic
(implemented in \texttt{sympy}) with a verified orthonormal basis for the
standard representation. \Cref{thm:level23} uses numerical computation
(\texttt{numpy}) over all $73810$ valid configurations.

The full verification suite consists of $12$ scripts,
publicly available at:
\begin{center}
\url{https://anonymous.4open.science/r/barrington-mixing}
\end{center}

\begin{center}
\begin{tabular}{clc}
\toprule
Script & Content & Method \\
\midrule
01 & Ambivalence of $S_5$ (similarity blindness) & Exhaustive \\
02 & Collision analysis, $\ell = 4$ & Exhaustive \\
03 & Collision scaling, $\ell = 4 \to 8$ & Sampled \\
04 & Fourier spectral analysis & Exact + sampled \\
05 & Barrington construction + Kilian obfuscation & Exhaustive \\
06 & Contraction operator (unconditional distribution) & Analytic \\
07 & T-independence proof verification & Exhaustive \\
08 & Conjunctions with wildcards & Exhaustive \\
09 & Universal mixing investigation & Exact + sampled \\
10 & Complete spectral contraction proof & Exact + numerical \\
11 & Eigenvector analysis (simplex discovery) & Numerical \\
12 & Level-$3$ exhaustive verification & Numerical \\
\bottomrule
\end{tabular}
\end{center}

Scripts 10--12 collectively verify the main theorem in under $15$ seconds
on a standard laptop (the exact-arithmetic steps in Script~10 dominate).
Scripts~01--09 provide additional cross-validation and extensions.


\bibliography{refs}

\end{document}
